\chapter{LDCS - Longest Decreasing Common Subsequence}

\section{Definizione formale del problema}

\subsection{Istanza}
Due sequenze finite $X = \langle x_1, x_2, \dots, x_m \rangle$ e $Y = \langle y_1, y_2, \dots, y_n \rangle$ di elementi appartenenti a un insieme $\Sigma$ munito di una relazione d'ordine totale stretta ($>$).

\subsection{Soluzione}
Una più lunga sottosequenza comune a $X$ e $Y$ ordinata in modo strettamente decrescente, ovvero una sequenza $Z = \langle z_1, z_2, \dots, z_k \rangle$ tale che:
\begin{itemize}
    \item $Z$ sia sottosequenza comune di $X$ e $Y$;
    \item $z_1 > z_2 > \dots > z_k$;
    \item la lunghezza $k = |Z|$ sia massima.
\end{itemize}

---

\section{Sottoproblema vincolato e coefficiente}

Per garantire la proprietà della sottostruttura ottima retrospettiva e verificare la decrescenza rispetto all'ultimo carattere comune scelto, si definisce il sottoproblema ausiliario vincolato $\text{LDCS}_V(X_i, Y_j)$.

\subsection{Sottoproblema vincolato $\text{LDCS}_V(X_i, Y_j)$}
Dati i prefissi $X_i = \langle x_1, \dots, x_i \rangle$ e $Y_j = \langle y_1, \dots, y_j \rangle$ per ogni $i \in \{0, \dots, m\}$ e $j \in \{0, \dots, n\}$, si cerca la più lunga sottosequenza comune decrescente \textbf{vincolata a terminare tassativamente con l'elemento comune} $x_i = y_j$. Se $x_i \neq y_j$, la sottosequenza vincolata non esiste (restituisce la sequenza vuota $\langle \rangle$).

\subsection{Coefficiente $c_{i,j}$}
Definiamo il coefficiente $c_{i,j}$ come la lunghezza della sottosequenza ottima vincolata:
\begin{equation*}
    c_{i,j} = |\text{LDCS}_V(X_i, Y_j)|
\end{equation*}

---

\section{Equazioni di ricorrenza}

\subsection{Casi base}
Se almeno uno dei due prefissi è vuoto ($i = 0 \lor j = 0$), oppure se i caratteri terminali non coincidono ($x_i \neq y_j$), non è possibile formare alcuna sottosequenza comune vincolata a terminare con $x_i = y_j$:
\begin{equation*}
    (i = 0 \lor j = 0 \lor x_i \neq y_j) \implies c_{i,j} = 0
\end{equation*}

\subsection{Passo ricorsivo}
Per $i > 0$, $j > 0$ e $x_i = y_j$, l'elemento comune $x_i$ viene concatenato alla migliore sottosequenza comune decrescente precedente terminante in una coppia di indici compatibile $x_h = y_k$ (con $1 \le h < i$ e $1 \le k < j$) che rispetti la decrescenza stretta ($x_h > x_i$):
\begin{equation*}
    c_{i,j} = \max \Big(\{ c_{h,k} : 1 \le h < i \land 1 \le k < j \land x_h = y_k \land x_h > x_i \} \cup \{0\}\Big) + 1
\end{equation*}
con la convenzione $\max \emptyset = 0$. Se nessun predecessore comune è maggiore di $x_i$, si ottiene $c_{i,j} = 0 + 1 = 1$, corrispondente alla sottosequenza formata dal solo elemento $\langle x_i \rangle$.

\subsection{Soluzione del problema principale}
Poiché la sottosequenza ottima globale non termina necessariamente con gli ultimi elementi delle sequenze ($x_m = y_n$), il valore ottimo globale del problema non vincolato corrisponde al massimo calcolato su tutta la matrice:
\begin{equation*}
    \text{Soluzione} = \max_{\substack{1 \le i \le m \\ 1 \le j \le n}} \{ c_{i,j} \}
\end{equation*}

---

\section{Algoritmo Bottom-Up}

La procedura iterativa compila la matrice delle lunghezze $c[0 \dots m, 0 \dots n]$ e la matrice dei puntatori ai predecessori $b[0 \dots m, 0 \dots n]$, mantenendo aggiornati ad ogni passo il massimo globale $\mathit{max\_tot}$ e le coordinate $(\mathit{ind\_m1}, \mathit{ind\_m2})$ della cella terminale ottima.

\begin{algorithm}[H]
\caption{Calcolo della tabella LDCS e dei predecessori}
\label{alg:ldcs-bottom-up}
\begin{algorithmic}[1]
\Procedure{LDCS}{$X, Y$}
    \State $m = \text{length}(X)$ \Comment{Dimensione della sequenza $X$}
    \State $n = \text{length}(Y)$ \Comment{Dimensione della sequenza $Y$}
    \State alloca matrice $c[0 \dots m, 0 \dots n]$ inizializzata a $0$
    \State alloca matrice $b[0 \dots m, 0 \dots n]$ inizializzata a $(0, 0)$
    \State $\mathit{max\_tot} = 0$ \Comment{Traccia della lunghezza ottima globale}
    \State $\mathit{ind\_m1} = 0$, $\mathit{ind\_m2} = 0$ \Comment{Coordinate della cella terminale dell'ottimo}
    \For{$i = 1 \textbf{ to } m$}
        \For{$j = 1 \textbf{ to } n$}
            \If{$x_i == y_j$} \Comment{Match comune: ricerca del miglior predecessore}
                \State $\mathit{tmp} = 0$ \Comment{Valore neutro per $\max \emptyset = 0$}
                \State $h_{\max} = 0$, $k_{\max} = 0$ \Comment{Puntatori neutri al predecessore}
                \For{$h = 1 \textbf{ to } i - 1$}
                    \For{$k = 1 \textbf{ to } j - 1$}
                        \If{$x_h == y_k \textbf{ and } x_h > x_i \textbf{ and } c[h, k] > \mathit{tmp}$}
                            \State $\mathit{tmp} = c[h, k]$ \Comment{Aggiorna il massimo tra i predecessori compatibili}
                            \State $h_{\max} = h$
                            \State $k_{\max} = k$
                        \EndIf
                    \EndFor
                \EndFor
                \State $c[i, j] = \mathit{tmp} + 1$ \Comment{Estende la sottosequenza ottima precedente}
                \State $b[i, j] = (h_{\max}, k_{\max})$ \Comment{Memorizza i puntatori di risalita}
                \If{$c[i, j] > \mathit{max\_tot}$} \Comment{Aggiornamento dell'ottimo globale}
                    \State $\mathit{max\_tot} = c[i, j]$
                    \State $\mathit{ind\_m1} = i$
                    \State $\mathit{ind\_m2} = j$
                \EndIf
            \Else
                \State $c[i, j] = 0$ \Comment{Caratteri diversi: nessuna LDCS vincolata valida}
                \State $b[i, j] = (0, 0)$
            \EndIf
        \EndFor
    \EndFor
    \State \Return $c, b, \mathit{max\_tot}, \mathit{ind\_m1}, \mathit{ind\_m2}$
\EndProcedure
\end{algorithmic}
\end{algorithm}

---

\section{Algoritmo di Ricostruzione}

La procedura risale a ritroso i puntatori memorizzati in $b$ partendo dalle coordinate dell'ottimo globale $(\mathit{ind\_m1}, \mathit{ind\_m2})$. La stampa posizionata al ritorno della chiamata ricorsiva (post-ordine) assicura l'emissione degli elementi nell'ordine originale da sinistra verso destra.

\begin{algorithm}[H]
\caption{Ricostruzione e stampa della LDCS ottima}
\label{alg:print-ldcs}
\begin{algorithmic}[1]
\Procedure{Print-LDCS}{$i, j, b, X$}
    \If{$i == 0 \lor j == 0$} \Comment{Terminatore di risalita}
        \State \Return
    \EndIf
    \State $(h, k) = b[i, j]$ \Comment{Estrae gli indici del predecessore}
    \State \Call{Print-LDCS}{$h, k, b, X$} \Comment{Chiamata ricorsiva a ritroso}
    \State $\text{print}(x_i)$ \Comment{Stampa in risalita l'elemento della soluzione}
\EndProcedure
\end{algorithmic}
\end{algorithm}

\subsection{Invocazione principale (Driver)}
Per avviare la ricostruzione sulla soluzione calcolata dall'algoritmo \textsc{LDCS}, si verifica preventivamente l'esistenza di una sottosequenza non vuota:
\begin{algorithmic}[1]
\If{$\mathit{max\_tot} == 0$}
    \State $\text{print}(\text{"Non esiste alcuna sottosequenza comune decrescente non vuota"})$
\Else
    \State \Call{Print-LDCS}{$\mathit{ind\_m1}, \mathit{ind\_m2}, b, X$}
\EndIf
\end{algorithmic}

---

\section{Analisi di Complessità}

\subsection{Complessità temporale}
\begin{itemize}
    \item \textbf{\textsc{LDCS} (Bottom-Up):} I due cicli esterni su $i$ e $j$ compiono $m \times n$ iterazioni. Per ogni cella con $x_i = y_j$, i due cicli annidati retrospettivi scandiscono $(i-1) \times (j-1)$ celle a costo unitario. Nel caso pessimo la complessità temporale è:
    \begin{equation*}
        T(m, n) = O\left( \sum_{i=1}^{m} \sum_{j=1}^{n} i \cdot j \right) = O(m^2 \cdot n^2)
    \end{equation*}
    \item \textbf{\textsc{Print-LDCS} (Ricostruzione):} A ciascun passo ricorsivo si retrocede di un match lungo la sequenza dei puntatori. Il numero di chiamate è pari alla lunghezza $k = \mathit{max\_tot}$ della sequenza ottima ($k \le \min(m, n)$), richiedendo tempo $O(\min(m, n))$.
\end{itemize}

\subsection{Complessità spaziale}
L'algoritmo alloca due matrici bidimensionali $c$ e $b$ di dimensioni $(m+1) \times (n+1)$, con costo in memoria ausiliaria pari a $\Theta(m \cdot n)$. Lo stack di ricorsione della procedura \textsc{Print-LDCS} occupa spazio pari a $O(\min(m, n))$.